\documentclass[10pt,a4paper]{amsart}
\usepackage[margin=19mm]{geometry}
\usepackage{amsmath,amssymb,mathtools}
\usepackage[scr=rsfs,cal=euler]{mathalfa}
\usepackage{xfrac}
\usepackage[unicode,hidelinks,bookmarks=false]{hyperref}
\newcommand{\ie}{i.e.\ }
\newcommand{\cf}{cf.\ }
\newcommand{\resp}{respectively\ }
\newcommand{\Subsection}{\subsection}
\providecommand{\nobreakdash}{\nobreak\hbox{-}}
\setlength{\emergencystretch}{2em}
\multlinegap0pt
\usepackage{cleveref}
\usepackage{aliascnt}
\usepackage{amsthm}
\usepackage{mathtools}
\usepackage{amsmath,amssymb,amsfonts}
\usepackage{graphicx}
\usepackage{bm}
\usepackage{enumitem}
\usepackage{mathrsfs}
\usepackage{xcolor}
\usepackage{tikz}
\usetikzlibrary{arrows.meta,calc,positioning,patterns,decorations.pathmorphing,decorations.markings}
\providecommand{\F}{\mathbb F}
\DeclareMathOperator{\rank}{rank}
\DeclareMathOperator{\row}{row}
\numberwithin{equation}{section}
\newtheorem{theorem}{Theorem}[section]
\newtheorem*{maintheorem}{Main Theorem}
\newaliascnt{lemma}{theorem}
\newtheorem{lemma}[lemma]{Lemma}
\aliascntresetthe{lemma}
\providecommand*{\lemmaautorefname}{Lemma}
\newaliascnt{proposition}{theorem}
\newtheorem{proposition}[proposition]{Proposition}
\aliascntresetthe{proposition}
\providecommand*{\propositionautorefname}{Proposition}
\newaliascnt{corollary}{theorem}
\newtheorem{corollary}[corollary]{Corollary}
\aliascntresetthe{corollary}
\providecommand*{\corollaryautorefname}{Corollary}
\newaliascnt{claim}{theorem}
\newtheorem{claim}[claim]{Claim}
\aliascntresetthe{claim}
\providecommand*{\claimautorefname}{Claim}
\newaliascnt{conjecture}{theorem}
\newtheorem{conjecture}[conjecture]{Conjecture}
\aliascntresetthe{conjecture}
\providecommand*{\conjectureautorefname}{Conjecture}
\theoremstyle{definition}
\newtheorem{fact}{Fact}[section]
\crefname{fact}{Fact}{Facts}
\Crefname{fact}{Fact}{Facts}
\newaliascnt{definition}{theorem}
\newtheorem{definition}[definition]{Definition}
\aliascntresetthe{definition}
\providecommand*{\definitionautorefname}{Definition}
\newaliascnt{example}{theorem}
\newtheorem{example}[example]{Example}
\aliascntresetthe{example}
\providecommand*{\exampleautorefname}{Example}
\newaliascnt{notation}{theorem}
\newtheorem{notation}[notation]{Notation}
\aliascntresetthe{notation}
\providecommand*{\notationautorefname}{Notation}
\newaliascnt{question}{theorem}
\newtheorem{question}[question]{Question}
\aliascntresetthe{question}
\providecommand*{\questionautorefname}{Question}
\theoremstyle{remark}
\newaliascnt{remark}{theorem}
\newtheorem{remark}[remark]{Remark}
\aliascntresetthe{remark}
\providecommand*{\remarkautorefname}{Remark}
\definecolor{tensorhighlight}{gray}{0.92}
\tikzset{
  graphvertex/.style={circle,fill,inner sep=1.35pt},
  leafvertex/.style={circle,draw=black,fill=black,inner sep=1.35pt},
  facevertex/.style={circle,draw,fill=white,inner sep=1.5pt,font=\scriptsize},
  graphlabel/.style={font=\scriptsize,inner sep=1pt}
}
\title[Represented tensor products of binary matroids]{Represented Tensor Products of Binary Matroids: the outerplanar cographic construction (Proposition 3.2)}
\author[H. Fu, Y. Ma, S. Wang]{Houshan Fu, Yujiao Ma, Suijie Wang}
\date{Source: arXiv:2609.29227v1 (CC BY 4.0)}
\begin{document}
\maketitle
\vspace{1em}
\noindent\textit{Source and licence.} Represented Tensor Products of Binary Matroids. Version arXiv:2609.29227v1 (CC BY 4.0), distributed under the Creative Commons Attribution 4.0 International licence, http://creativecommons.org/licenses/by/4.0/, as recorded in the arXiv OAI licence metadata for this exact version and shown on the abstract page https://arxiv.org/abs/2609.29227v1. This item is an editorial re-presentation of the paper's cographic construction for two-connected outerplanar graph factors, proving that the represented tensor product of a cycle matroid with an outerplanar graphic matroid is the dual of the cycle matroid of an explicitly constructed graph, printed as Proposition 3.2. A full rights notice is delivered at licenses/arxiv-2609.29227-CC-BY-4.0.txt.
\noindent\textit{Editorial scope.} Counted unit: the paper's cographic construction for two-connected outerplanar graph factors, proving that the represented tensor product of a cycle matroid with an outerplanar graphic matroid is the dual of the cycle matroid of an explicitly constructed graph, printed as Proposition 3.2 (source0.tex lines the statement at lines 1123--1133 and the proof at lines 1134--1304, with two short context spans that the construction refers to), delivered together with the paper's complete original proof of it as the single primary proof body. No mathematics is written, repaired or re-derived: statement and proof are the paper's own text line for line, in the paper's own notation. Required context retained uncounted: the two short spans of the preceding construction that the counted proof refers to. Every cross-reference inside the retained spans resolves inside the delivered document. Printed numbers are the paper's own: the wrapper restores the paper's counters before the retained material, so the counted statement prints as Proposition 3.2, the paper's own number for it, and the numbered display of the statement carries the paper's own equation number (3.2). Omitted from this package and not redistributed: the paper's preamble and front matter as such (of which only the title and the author names are reproduced in the wrapper); the sections, lemmas, theorems and examples that lie outside the retained spans; and the paper's bibliography with its bib data file. No citation command occurs in any retained span, so the delivered document carries no bibliography and no reference or citation command of the delivered text was rewritten or adapted. Printed slips kept literal, among them the paper's own wording and its own choice of symbols. Layout and class: the delivered wrapper is the lane's pinned amsart editorial wrapper at 10pt on A4, to which this package adds the paper's own theorem-like declarations and the shorthand macros its retained text uses, all copied verbatim from the paper's source. No publisher class file, no publisher style file and no upstream script is redistributed. Assets: no retained span of this package contains a figure environment or a graphics-inclusion command, no image or artwork file exists in or is redistributed with this package, and the manifest and the delivery evidence both carry an empty asset-dependency list. A full rights notice is delivered at licenses/arxiv-2609.29227-CC-BY-4.0.txt.
\setcounter{section}{2}
\begin{equation}\label{eq:cycle-space-dual-representation}
 A\text{ has full row rank and }\row(A)=Z_{\F}(G)
 \quad\Longrightarrow\quad
 M[A]\cong M^*(G).
\end{equation}

\setcounter{section}{3}
\setcounter{theorem}{4}
\setcounter{equation}{3}
\begin{equation}\label{eq:vertex-edge}
V(X_{m,G})=([m]\times\mathcal F(G))\sqcup\{x_e:e\in E_\partial(G)\},\qquad E(X_{m,G})=[m]\times E(G).
\end{equation}

\begin{proposition}
\label{prop:outerplanar-cographic}
Let \(m\ge3\), \(\F\) be a field, and \(G\) be a \(2\)-connected outerplanar simple graph.  Then the graph \(X_{m,G}\) constructed above satisfies
\[
M_{\F}(C_m,G)\cong M^*(X_{m,G}).
\]
Moreover, if \(G=C_n\), where \(n\ge3\), then
\begin{equation}\label{eq:two-cycle-identity}
 M_{\F}(C_m,C_n)\cong M^*(K_{m,n}).
\end{equation}
\end{proposition}

\begin{proof}
Set \(A_{C_m}=\Delta_m\), whose \(i\)-th row is
\(\boldsymbol e_i^{(m)}-\boldsymbol e_m^{(m)}\), and let \(A_G\)
be the reduced incidence matrix obtained by deleting the row at
\(v_0\in V(G)\).  These are full row rank representations of the
factors; the rows of \(A_G\) are
\(\boldsymbol b_G(v)\) for \(v\ne v_0\).

For each such \(v\), the graph \(G-v\) is connected, so
\(\partial_G(v)\) is a bond.  Its dual edges form a cycle
\(C_v^*\) of \(G^*\).  The two outer edges incident with \(v\)
correspond to the two edges of \(C_v^*\) incident with
\(v_\infty\).  Splitting \(v_\infty\) therefore turns \(C_v^*\)
into a path \(P_v\) between the corresponding leaves of
\(\widehat T_G\).

Traverse \(C_v^*\) with the face of \(G^*\) corresponding to
\(v\) on the right, and give \(P_v\) the induced orientation.
At an edge directed away from \(v\), this traversal agrees with
the dual-edge orientation; at an edge directed towards \(v\),
it is opposite.  Thus the signed edge vector of \(P_v\) is
\(\boldsymbol b_G(v)\).  Let \(P_v^j\) be its copy in layer
\(j\in[m]\).  For \(i\in[m-1]\), the paths \(P_v^i\) and
\(P_v^m\) have the same endpoints and disjoint interiors, so
their union is a cycle \(\Gamma_{i,v}\).  Traverse \(P_v^i\)
forward and \(P_v^m\) backward.  The resulting cycle vector is
\[
 \boldsymbol{z}_{\Gamma_{i,v}}=(\boldsymbol{e}_i-\boldsymbol{e}_m) \otimes\boldsymbol b_G(v).
\]
Figure~\ref{fig:dual-path-cycle} illustrates the construction of these cycle vectors.
\begin{figure}[!ht]
\centering
\begin{tikzpicture}[
  x=.64cm,y=.75cm,
  line cap=round,
  line join=round,
  line width=.60pt,
  >=Stealth,
  localface/.style={circle,draw=black,fill=blue!9,minimum size=7.5mm,inner sep=1pt,font=\scriptsize},
  localouter/.style={circle,draw=black,fill=black!7,minimum size=8.5mm,inner sep=1pt,font=\scriptsize},
  localleaf/.style={circle,fill=black,inner sep=1.35pt},
  ilayer/.style={circle,draw=black,fill=blue!11,minimum size=10.5mm,inner sep=0pt,font=\scriptsize},
  mlayer/.style={circle,draw=black,fill=orange!17,minimum size=10.5mm,inner sep=0pt,font=\scriptsize},
  locallab/.style={font=\scriptsize,inner sep=1pt},
  localtitle/.style={font=\scriptsize\bfseries},
  patharrow/.style={
    postaction={decorate},
    decoration={markings,mark=at position .55 with
      {\arrow{Stealth[length=1.8mm,width=1.2mm]}}}
  }
]
  \node[localtitle,anchor=center] at (4.10,1.55)
    {The dual cycle \(C_v^*\)};
  \node[localouter] (vinf) at (0,0) {\(v_\infty\)};
  \node[localface] (c1) at (2.0,.75) {\(F_1\)};
  \node[localface] (c2) at (4.0,.75) {\(F_2\)};
  \node[localface] (cdm2) at (6.55,.75) {\(F_{d-2}\)};
  \node[localface] (cdm1) at (8.65,.75) {\(F_{d-1}\)};
  \draw (vinf.north east)--node[locallab,above left=2pt]{\(e_1^*\)}(c1.west);
  \draw (c1.east)--node[locallab,above=2pt]{\(e_2^*\)}(c2.west);
  \coordinate (cLeft) at ($(c2.east)+(.34,0)$);
  \coordinate (cRight) at ($(cdm2.west)+(-.34,0)$);
  \draw (c2.east)--(cLeft);
  \foreach \t in {.25,.50,.75}
    \fill ($(cLeft)!\t!(cRight)$) circle[radius=.65pt];
  \draw (cRight)--(cdm2.west);
  \draw (cdm2.east)--node[locallab,above=2pt]{\(e_{d-1}^*\)}(cdm1.west);
  \draw (cdm1.south west) to[bend left=23]
        node[locallab,below=6pt]{\(e_d^*\)}(vinf.south east);

  \draw[-{Stealth},line width=.65pt] (10.0,.35)--(11.2,.35)
        node[midway,above=5pt,locallab]{split \(v_\infty\)};

  \node[localtitle,anchor=center] at (16.73,1.55)
    {Splitting \(v_\infty\) produces the oriented path \(P_v\)};
  \node[localleaf] (x1) at (12.1,.35) {};
  \node[localface] (p1) at (13.75,.35) {\(F_1\)};
  \node[localface] (p2) at (15.45,.35) {\(F_2\)};
  \node[localface] (pdm2) at (17.85,.35) {\(F_{d-2}\)};
  \node[localface] (pdm1) at (19.75,.35) {\(F_{d-1}\)};
  \node[localleaf] (xd) at (21.35,.35) {};
  \node[locallab,anchor=north east] at ($(x1)+(0,-.28)$) {\(x_{e_1}\)};
  \node[locallab,anchor=north west] at ($(xd)+(0,-.28)$) {\(x_{e_d}\)};
  \draw[patharrow] (x1.center)--node[locallab,above=4pt]{\(e_1^*\)}(p1.west);
  \draw[patharrow] (p1.east)--node[locallab,above=4pt]{\(e_2^*\)}(p2.west);
  \coordinate (pLeft) at ($(p2.east)+(.30,0)$);
  \coordinate (pRight) at ($(pdm2.west)+(-.30,0)$);
  \draw (p2.east)--(pLeft);
  \foreach \t in {.25,.50,.75}
    \fill ($(pLeft)!\t!(pRight)$) circle[radius=.65pt];
  \draw (pRight)--(pdm2.west);
  \draw[patharrow] (pdm2.east)--node[locallab,above=4pt]{\(e_{d-1}^*\)}(pdm1.west);
  \draw[patharrow] (pdm1.east)--node[locallab,above=4pt]{\(e_d^*\)}(xd.center);

  \draw[-{Stealth},line width=.65pt] (10.65,-1.25)--(10.65,-2.85)
        node[midway,right=3pt,locallab]{take layers \(i\) and \(m\)};

  \node[localtitle,anchor=center] at (10.68,-3.05)
    {The two layer copies form \(\Gamma_{i,v}\)};
  \node[localleaf,label=left:{\scriptsize \(x_{e_1}\)}] (y1) at (2.1,-5.35) {};
  \node[localleaf,label=right:{\scriptsize \(x_{e_d}\)}] (yd) at (19.25,-5.35) {};
  \node[ilayer] (i1) at (5.0,-4.35) {\(i,F_1\)};
  \node[ilayer] (i2) at (7.6,-4.35) {\(i,F_2\)};
  \node[ilayer] (idm2) at (13.65,-4.35) {\(i,F_{d-2}\)};
  \node[ilayer] (idm1) at (16.35,-4.35) {\(i,F_{d-1}\)};
  \node[mlayer] (m1) at (5.0,-6.35) {\(m,F_1\)};
  \node[mlayer] (m2) at (7.6,-6.35) {\(m,F_2\)};
  \node[mlayer] (mdm2) at (13.65,-6.35) {\(m,F_{d-2}\)};
  \node[mlayer] (mdm1) at (16.35,-6.35) {\(m,F_{d-1}\)};

  \draw[patharrow] (y1.center)--(i1.west);
  \draw[patharrow] (i1.east)--(i2.west);
  \coordinate (iLeft) at ($(i2.east)+(1.55,0)$);
  \coordinate (iRight) at ($(idm2.west)+(-1.55,0)$);
  \draw[patharrow] (i2.east)--(iLeft);
  \foreach \t in {.25,.50,.75}
    \fill ($(iLeft)!\t!(iRight)$) circle[radius=.65pt];
  \draw[patharrow] (iRight)--(idm2.west);
  \draw[patharrow] (idm2.east)--(idm1.west);
  \draw[patharrow] (idm1.east)--(yd.center);
  \draw[patharrow] (yd.center)--(mdm1.east);
  \draw[patharrow] (mdm1.west)--(mdm2.east);
  \coordinate (mLeft) at ($(m2.east)+(1.55,0)$);
  \coordinate (mRight) at ($(mdm2.west)+(-1.55,0)$);
  \draw[patharrow] (mdm2.west)--(mRight);
  \foreach \t in {.25,.50,.75}
    \fill ($(mLeft)!\t!(mRight)$) circle[radius=.65pt];
  \draw[patharrow] (mLeft)--(m2.east);
  \draw[patharrow] (m2.west)--(m1.east);
  \draw[patharrow] (m1.west)--(y1.center);
  \node[locallab,text=blue!55!black] at (10.65,-4.88) {\(P_v^i\)};
  \node[locallab,text=orange!70!black] at (10.65,-5.82) {\(P_v^m\)};
  \node[font=\footnotesize] at (10.65,-7.88)
    {\(\displaystyle
      \Gamma_{i,v}=P_v^i\cup P_v^m,
      \qquad
      \boldsymbol{z}_{\Gamma_{i,v}}
      =(\boldsymbol{e}_i-\boldsymbol{e}_m)\otimes
        \boldsymbol b_G(v).\)};
\end{tikzpicture}
\caption{Splitting \(v_\infty\) turns \(C_v^*\) into \(P_v\);
the oppositely traversed copies in layers \(i\) and \(m\) form
\(\Gamma_{i,v}\).  Here \(d=\deg_G(v)\), and the incident edges
\(e_1,\ldots,e_d\) are indexed in the order in which their dual
edges are traversed along \(C_v^*\), starting at \(v_\infty\).
The edges \(e_1\) and \(e_d\) are the two outer edges.}
\label{fig:dual-path-cycle}
\end{figure}
The vectors \(\boldsymbol z_{\Gamma_{i,v}}\) are precisely the rows of \(A_{C_m}\otimes A_G\).  Consequently,
\begin{equation}\label{eq:outerplanar-rowspace-inclusion}
 \row(A_{C_m}\otimes A_G)\subseteq Z_{\F}(X_{m,G}).
\end{equation}

By Euler's formula and \eqref{eq:vertex-edge},
\[
 \begin{aligned}
 \dim Z_{\F}(X_{m,G})
 &=m|E(G)|-|V(X_{m,G})|+1\\
 &=(m-1)(|V(G)|-1)
 =\rank(A_{C_m}\otimes A_G).
 \end{aligned}
\]
Thus \eqref{eq:outerplanar-rowspace-inclusion} is an equality, and
\eqref{eq:cycle-space-dual-representation} gives
\(M_{\F}(C_m,G)\cong M^*(X_{m,G})\).

If \(G=C_n\), then \(\widehat T_G\cong K_{1,n}\).
Identifying equally labelled leaves in \(m\) copies gives
\(X_{m,C_n}\cong K_{m,n}\), proving
\eqref{eq:two-cycle-identity}.
\end{proof}

\end{document}
